% Generated draft: requires mathematical and source review.
\documentclass[11pt]{article}
\usepackage[margin=25mm]{geometry}
\usepackage{fontspec}
\setmainfont{Latin Modern Roman}
\usepackage{amsmath,amssymb,mathtools,unicode-math}
\setmathfont{Latin Modern Math}
\usepackage{enumitem,needspace,xurl,hyperref}
\setlength{\parindent}{0pt}
\setlength{\parskip}{5pt}
\setlength{\emergencystretch}{3em}
\newcommand{\N}{\mathbb{N}}
\newcommand{\Z}{\mathbb{Z}}
\newcommand{\Q}{\mathbb{Q}}
\newcommand{\R}{\mathbb{R}}
\newcommand{\C}{\mathbb{C}}
\newcommand{\F}{\mathbb{F}}
\newcommand{\A}{\mathbb{A}}
\newcommand{\PP}{\mathbb{P}}
\newcommand{\E}{\mathbb{E}}
\newcommand{\eps}{\varepsilon}
\newcommand{\dd}{\,\mathrm{d}}
\begin{document}
\section*{UnsolvedMath ID: 701366}
\textbf{Source problem code:} catalogue entry
\subsection{Definitions and mathematical statement}
Let \(\mathcal P\) be the set of primes. Under the customary interpretation, the Nottingham group \(J=N(p)\) is the group
\[
N(p)=\left\{t+\sum_{k=2}^{\infty}a_k t^k:a_k\in\F_p\right\}
\]
under formal composition, with the topology in which congruence modulo \(t^m\), for increasing \(m\), defines neighborhoods of the identity. Here \(\F_p\) is the field with \(p\) elements. The supplied source refers to Archive Problem 12.24 for the definition of \(N(p)\), but does not reproduce that definition; agreement with this customary interpretation requires review.

A pro-\(p\) group is a topological group that is an inverse limit of finite \(p\)-groups. For a pro-\(p\) group \(H\), use the closed lower central series
\[
\gamma_1(H)=H,\qquad
\gamma_{i+1}(H)=\overline{\langle [x,y]:x\in\gamma_i(H),\ y\in H\rangle}
\quad(i\geq1),
\]
where \([x,y]=x^{-1}y^{-1}xy\), and the bar denotes topological closure. For a positive integer \(n\), call \(G\) an \(n\)-approximation to \(J=N(p)\) if \(G\) is an infinite pro-\(p\) group and
\[
G/\gamma_n(G)\cong J/\gamma_n(J).
\]
Here isomorphisms are interpreted as topological group isomorphisms. The source does not explicitly specify closure in the lower central series or the category of isomorphisms, so these conventions also require confirmation.

\textbf{Question (C. R. Leedham-Green, 2010).} Does there exist a function
\[
f:\mathcal P\longrightarrow\Z_{\geq1}
\]
such that, for every prime \(p\) and every infinite pro-\(p\) group \(G\),
\[
G/\gamma_{f(p)}(G)\cong N(p)/\gamma_{f(p)}(N(p))
\quad\Longrightarrow\quad
\gamma_i(G)/\gamma_{i+1}(G)
\cong
\gamma_i(N(p))/\gamma_{i+1}(N(p))
\quad\text{for every integer }i\geq1?
\]
The conclusion asks for isomorphism of each corresponding lower central factor; it does not prescribe compatible choices of these isomorphisms.

The source also refers to Problem 14.56 and states that an affirmative answer would imply the fact it describes as already known: \(N(p)\) is finitely presented as a pro-\(p\) group for \(p>2\); it refers to Problem 14.55 for this remark. This is a supplied background qualification, not an independently verified result here. The condition \(p>2\) qualifies that finite-presentation remark, whereas the supplied question states no such restriction on \(p\).
\subsection{Short English statement}
Does a finite initial lower central quotient of the Nottingham group, with its depth depending only on the prime \(p\), force every infinite pro-\(p\) group sharing that quotient to have all the same lower central factors? The supplied registry labels the problem open. The referenced definition of the Nottingham group and the topological conventions require confirmation.
\subsection{Sources}
\begin{itemize}
\item Local UnsolvedMath registry, permanent ID \(701366\), exact source code \texttt{catalogue entry}, title \emph{Kourovka Notebook Problem 17.64}. The supplied record labels the problem \texttt{open}; this label is not an independent assessment of its current mathematical status. The record was updated on 2026-10-16 and identifies upstream revision \texttt{d371ee917ac9647fbb06cd6fc529deec6f3dda82}.
\item \emph{The Kourovka Notebook (Unsolved Problems in Group Theory)}, original Problem 17.64, 17th issue (2010), proposed by C. R. Leedham-Green. The supplied background also identifies the 21st issue edition (2026), arXiv:1401.0300v48. The collection description names E. I. Khukhro and V. D. Mazurov as editors. These bibliographic details are supplied attribution and have not been independently checked.
\item Internal references supplied with the statement: Archive Problem 12.24 for the definition of \(N(p)\); Problem 14.56 for comparison; and Problem 14.55 for the finite-presentation background remark for \(p>2\). The contents of these referenced entries were not supplied.
\end{itemize}
\end{document}
